\documentclass{ximera}

\begin{document}

    \title{Lecture 10}
    
    \begin{abstract}  
    Outline: \\
    - Algebraic Limit Theorem for series\\
    - Cauchy Criterion for series\\
    - Comparison Test \\
    - Geometric series \\
    - Absolute Convergence Test \& Alternating Series Test
\end{abstract}

\maketitle

% \textcolor{red}{Go over the last problem on Quiz 1}

Remember that we can treat a series as a sequence of partial sums:
\begin{align*}
    \sum\limits_{k=1}^\infty a_k &= \lim\limits_{n\to\infty} s_n, \text{ where}
    \\
    s_n &= a_1 + a_2 + \cdots + a_n.
\end{align*}

\begin{exercise}
    Formulate parts (a) and (b) of the Algebraic Limit Theorem but now for infinite series.  Prove these statements are true.
\end{exercise}

\begin{exercise}
    What does the Algebraic Limit Theorem (c) and (d) say about the convergence of infinite series?  Where do the suppositions
    \begin{enumerate}
        \item $\sum\limits_{k=1}^\infty a_kb_k= AB$
        \item $\sum\limits_{k=1}^\infty a_k/\sum\limits_{k=1}^\infty b_k = A/B$ if $B\neq 0$
    \end{enumerate}
    go wrong?
\end{exercise}

\begin{theorem}{(Algebraic Limit Theorem for Series)}
    If $\sum\limits_{k=1}^\infty a_k = A$, $\sum\limits_{k=1}^\infty b_k =B$, then 
    \begin{enumerate}
        \item $\sum\limits_{k=1}^\infty ca_k = cA$ for all $c\in\mathbb{R}$
        \item $\sum\limits_{k=1}^\infty a_k +b_k = A+B$
    \end{enumerate}
\end{theorem}

% \textcolor{red}{Have them formulate the CC for series with you using the theorem that a sequence in $\mathbb{R}$ converges if and only if it is Cauchy.}

\begin{theorem}{(Cauchy Criterion for Series)}
    The series $\sum\limits_{k=1}^\infty a_k$ converges if and only if for all $\epsilon > 0$, there exists $N \in \mathbb{N}$ such that for all $n\geq m\geq N$, 
    \begin{align*}
        |a_{m+1} + a_{m+2} + \cdots + a_n| < \epsilon.
    \end{align*}
\end{theorem}
\begin{proof}
    $(\implies)$  Suppose the series converges.  Then for all $\epsilon > 0$, there exists an $N\in\mathbb{N}$ such that for all $n\geq m \geq N$ implies
    \begin{align*}
        |s_n - s_m| < \epsilon,
    \end{align*}
    which is exactly 
    \begin{align*}
        |a_{m+1} + a_{m+2} + \cdots + a_n| < \epsilon.
    \end{align*}

    $(\impliedby)$ Since for all $\epsilon > 0$ there exists  $N \in \mathbb{N}$ such that for all $n\geq m \geq N$, 
    \begin{align*}
        |a_{m+1} + a_{m+2} + \cdots + a_n| < \epsilon,
    \end{align*}
    but again, 
    \begin{align*}
        |a_{m+1} + a_{m+2} + \cdots + a_n| = |s_n - s_m|.
    \end{align*}
\end{proof}

This simple theorem leads to some other, very useful theorems.

\begin{theorem}
    If the series $\sum\limits_{k=1}^\infty a_k$ converges, then $\{a_k\} \to 0$.
\end{theorem}
\begin{proof}
    \textcolor{red}{How can we prove this using the Cauchy Criterion?}
    Let $\epsilon > 0$.  Fix $N$ such that for all $n,m \geq N$
    \begin{align*}
        |a_{m+1} + a_{m+2} + \cdots + a_n| < \epsilon.
    \end{align*}
    Choose $m = n-1$ (and therefore $n \geq N+1$).  Then $|a_n| < \epsilon$, or $|a_n - 0| < \epsilon$, and the proof is concluded.
\end{proof}

\begin{theorem}{(Comparison Test)}
    Assume $\{a_k\}, \{b_k\}$ are sequences satisfying $0 \leq a_k \leq b_k$ for all $k \in \mathbb{N}$.
    \begin{enumerate}
        \item If $\sum\limits_{k=1}^\infty b_k$ converges, then $\sum\limits_{k=1}^\infty a_k$ converges.
        \item If $\sum\limits_{k=1}^\infty a_k$ diverges, then $\sum\limits_{k=1}^\infty b_k$ diverges.
    \end{enumerate}
\end{theorem}

\begin{exercise}
    Prove the above theorem. \textcolor{red}{Prove using CC, and bonus: prove using Monotone Convergence Theorem.}
\end{exercise}

Geometric series: 
\begin{align*}
    \sum\limits_{k=0}^\infty a r^k = a + ar + ar^2 + ar^3 + \cdots
\end{align*}
For $r \neq 1$, there is an algebraic identity:
\begin{align*}
    (1-r)(1+r+r^2 + r^3 + \cdots + r^{m-1}) = 1-r^m
\end{align*}
\textcolor{red}{How can I use this identity to rewrite the partial sum?}

\begin{align}
    s_m = \frac{a(1-r^m)}{1-r}.
\end{align}
\textcolor{red}{When does this sequence of partial sums converge?}  Apply Algebraic Limit theorem for sequences multiply times (noting that $r^m$ converges if and only if $|r|< 1$, which we proved earlier in the course).

On your own, prove the following: 
\begin{theorem}{(Absolute Convergence Test)}
If $\sum\limits_{k=1}^\infty |a_k|$ converges, then $\sum\limits_{k=1}^\infty a_k$ converges.
\end{theorem}
and 
\begin{theorem}{(Alternating Series Test)}
 Let $\{a_k\}$ be a sequence such that 
 \begin{enumerate}
     \item $a_k \geq a_{k+1}$ for all $k \in \mathbb{N}$
     \item $a_k \to 0$.
 \end{enumerate}
Then the alternating series $\sum\limits_{k=1}^\infty (-1)^{k+1}a_k$ converges.
\end{theorem}

\begin{definition}
    If $\sum\limits_{k=1}^\infty |a_k|$ converges, then we say the original series $\sum\limits_{k=1}^\infty a_k$ converges absolutely.  Otherwise, if $\sum\limits_{k=1}^\infty a_k$ converges but $\sum\limits_{k=1}^\infty |a_k|$ does not, we say the $\sum\limits_{k=1}^\infty a_k$ converges conditionally. 
\end{definition}
\textcolor{red}{Given this definition, what can we say about the series $\sum\limits_{k=1}^\infty \frac1k$?}

\end{document}